\section{Examples}\label{sec:examples}

\subsection{Maximisers of the sum of distances}
\label{sec:maxim-sum-dist}
As pointed out in Remark~\ref{rem:invariance} Skriganov
\cite{Skriganov2020:stolarskys_invariance_principle} generalised Stolarsky's
invariance principle \cite{Stolarsky1973:sums_distances_sphere} to the
projective spaces $\mathbb{FP}^{d-1}$. As a consequence the $L^2$-discrepancy
is minimised amongst all sets of the same cardinality by configurations
maximising the sum of mutual ``chordal'' distances
\begin{equation}\label{eq:sum-dist}
  \sum_{x,y\in X_N}\sin(\theta(x,y)).
\end{equation}
It was shown in \cite{Skriganov2019:point_distributions_homogeneous} that there
exist positive constants $K_d$ and $k_d$ such that
\begin{equation}
  \label{eq:skriganov}
  \begin{split}
  -&K_d(\#X_N)^{1-\frac1D}\leq\\
  &\sum_{x,y\in X_N}\sin(\theta(x,y))-(\#X_N)^2\iint\limits_{(\mathbb{FP}^{d-1})^2}
  \sin(\theta(x,y))\,d\sigma(x)\,d\sigma(y)\\
  &\leq-k_d(\#X_N)^{1-\frac1D}
  \end{split}
\end{equation}
for a point set $X_N$ maximising \eqref{eq:sum-dist}.

Integrating \eqref{eq:var-Jacobi} against $\sin(2r)\,dr$ and using
\begin{align*}
  &\int_0^{\frac\pi2}\sin(r)^{4\alpha+4}\cos(r)^{4\beta+4}
  P_{n-1}^{(\alpha+1,\beta+1)}(\cos(2r))^2\sin(2r)\,dr=\\
  &\frac{\Gamma(2\alpha+3)\Gamma(2\beta+3)}{\alpha+\beta+2}
  2^{2n-2}(n+\alpha+\beta+1)\times\\
  &\frac{(\frac12)_{n-1}(\alpha+2)_{n-1}(\beta+2)_{n-1}(\alpha+\beta+2)_{n-1}}
  {((n-1)!)^2\Gamma(2(n+\alpha+\beta+3))}
\end{align*}
(see \cite{Skriganov2020:stolarskys_invariance_principle})
we obtain
\begin{equation}
  \label{eq:l2-discr}
  \begin{split}
    &\mathrm{D}_{L^2}(X_N)^2=C_{\alpha,\beta}^2\Gamma(2\alpha+3)\Gamma(2\beta+3)
    \times\\
    &\sum_{n=1}^\infty 2^{2n-2}(2n+\alpha+\beta+1)
    \frac{(\frac12)_{n-1}(\alpha+\beta+1)_n(\alpha+\beta+2)_n}
    {(\alpha+1)_n\Gamma(2(n+\alpha+\beta+3))}\times\\
    &\sum_{x,y\in X_N}
    P_n^{(\alpha,\beta)}(\cos(2\theta(x,y)))
  \end{split}
\end{equation}
(see also \cite{Skriganov2020:stolarskys_invariance_principle,
  Skriganov2019:stolarskys_invariance_principle_ii}).
We notice that the coefficients in the sum are of the order $n^{-\alpha-2}$.

We will use the inequality
\begin{equation}\label{eq:jacobi-bound}
  \begin{split}
    &\sin(\theta)^{a+\frac12}\cos(\theta)^{b+\frac12}
    |P_n^{(a,b)}(\cos(2\theta))|\\
    &\leq\frac{\Gamma(a+1)}{\sqrt\pi}
    \binom{n+a}n
    \left(n+\frac12(a+b+1)\right)^{-a-\frac12}
    =\mathcal{O}\left( n^{-\frac12}\right)
  \end{split}
\end{equation}
valid for $a\geq b>-1$ (see
\cite{Chow_Gatteschi_Wong1994:bernstein_type_inequality}). Inserting
$a=\alpha+1$ and $b=\beta+1$ in \eqref{eq:jacobi-bound} we estimate the number
variance by
\begin{equation}\label{eq:variance-upper}
  \mathbb{V}(X_N,r)\leq C\sin(r)^{2\alpha+1}\cos(r)^{2\beta+1}
  \sum_{n=1}^\infty n^{-\alpha-2}
  \sum_{x,y\in X_N}P_n^{(\alpha,\beta)}(\cos(2\theta(x,y))).
\end{equation}
Putting \eqref{eq:l2-discr} and~\eqref{eq:variance-upper} together and
observing that $2\alpha+1=D-1$ we obtain
\begin{equation}
  \label{eq:hyper}
  \mathbb{V}(X_N,r)=\mathcal{O}\left( \sin(r)^{D-1}\mathrm{D}_{L^2}(X_N)\right).
\end{equation}
For point sets with optimal $L^2$-discrepancy (or equivalently, maximal sum of
distances) we obtain
\begin{equation*}
  \mathbb{V}(X_N,r)=\mathcal{O}\left( \sin(r)^{D-1}(\#X_N)^{1-\frac1D}\right).
\end{equation*}
This inequality immediately implies the theorem.
\begin{theorem}
  Let $(X_N)_{N\in\mathbb{N}}$ be a sequence of maximisers of the sum of mutual
  distances \eqref{eq:sum-dist}. Then this sequence is hyperuniform in all all
  three regimes.
\end{theorem}
\subsection{Jittered sampling}\label{sec:jittered-sampling}
A very simple and obvious probabilistic model for a point set $X_N$ is given by
jittered sampling. We start with a partition of the space $\mathbb{FP}^{d-1}$
into $N$ sets $A_j$ ($j=1,\ldots,N$) of equal area and diameter
$\leq CN^{-\frac1D}$, where the constant only depends on the space. It is known
from \cite{Gigante_Leopardi2017:diameter_ahlfors} that such a partition exists.

We define a point process $\mathscr{X}_N^J$, called the \emph{jittered sampling
  process}, by choosing $N$ points $x_j\in A_j$ ($j=1,\ldots,N$), where each
point is chosen with respect to the surface measure restricted to $A_j$. The
number variance of this process is then given by
\begin{equation}\label{eq:variance-jittered}
  \begin{split}
    &V_N=\mathbb{V}(\mathscr{X}_N^J,r)\\
    &=\!\!\! \int\limits_{\mathbb{FP}^{d-1}}\int\limits_{A_1}\!\!\cdots\!\!\int\limits_{A_N}
    \!\!\left(\sum_{j=1}^N\mathbbm{1}_{B(x,r)}(x_j)-N\sigma(B(x,r))\right)^2\!\!
    d\sigma_1(x_1)\cdots d\sigma_N(x_N)\,d\sigma(x),
  \end{split}
\end{equation}
where $\sigma_j(A)=N\sigma(A_j\cap A)$ denotes the surface measure restricted
to $A_j$.

Expanding the square in~\eqref{eq:variance-jittered} and using the fact that
$\mathbbm{1}_{B(y,r)}(x)=\mathbbm{1}_{B(x,r)}(y)$ we obtain
\begin{align*}
  V_N&=\sum_{i\neq j}\int_{A_i}\int_{A_j}
  \sigma(B(x_i,r)\cap B(x_j,r))\,d\sigma_i(x_i)\,
  d\sigma_j(x_j)\\
  &+N\sigma(B(\cdot,r))-N^2\sigma(B(\cdot,r))^2 \\
  &=N^2\int_{\mathbb{FP}^{d-1}}\int_{\mathbb{FP}^{d-1}}
  \sigma(B(x,r)\cap B(y,r))\,d\sigma(x)\,
  d\sigma(y)\\
  &-
  \sum_{i=1}^N\int_{A_i}\int_{A_i}\sigma(B(x_i,r)\cap B(x_j,r))\,d\sigma_i(x_i)\,
  d\sigma_j(x_j)\\
  &+N\sigma(B(\cdot,r))-N^2\sigma(B(\cdot,r))^2\\
  &=\frac12\sum_{i=1}^N\int_{A_i}\int_{A_i}\sigma(B(x_i,r)\triangle B(y_i,r))
  \,d\sigma_i(x_i)\,d\sigma_i(y_i),
\end{align*}
where we have used
\begin{equation*}
  \int_{\mathbb{FP}^{d-1}}\int_{\mathbb{FP}^{d-1}}
  \sigma(B(x,r)\cap B(y,r))\,d\sigma(x)\,
  d\sigma(y)=\sigma(B(\cdot,r))^2
\end{equation*}
and used the notation $A\triangle B=A\cup B\setminus(A\cap B)$ for the
symmetric difference. A similar computation has been used in
\cite{Brauchart_Grabner_Kusner+2020:hyperuniform_point_sets}.

The measure of the symmetric difference of two balls can be bounded
\begin{equation}\label{eq:symmetric}
  \sigma(B(x,r)\triangle B(y,r))\leq
  C'\theta(x,y)\mathrm{surface}(\partial B(\cdot,r))
\end{equation}
for some constant $C'>0$. Since the diameter of each part $A_j$ of the partition
is bounded by $CN^{-\frac1D}$ we obtain the bound
\begin{equation*}
  V_N\leq CC'N^{1-\frac1D}\mathrm{surface}(\partial B(\cdot,r))=
  \mathcal{O}(r^{D-1}N^{1-\frac1D}),
\end{equation*}
from which we derive the theorem.
\begin{theorem}
  The jittered sampling process $\mathscr{X}_N^J$ is hyperuniform in all
three regimes.
\end{theorem}
\subsection{The harmonic ensemble}\label{sec:harmonic-ensemble}
Determinantal point processes (see
\cite{Hough_Krishnapur_Peres+2009:zeros_gaussian}) have become a very useful
tool to provide probabilistic models, which behave better than i.i.d. models
due to built mutual repulsion of the points. Especially, in the context of
finding well distributed point sets such processes have proven to be a very
useful tool to establish good bounds for various quality measures of point
distributions (see, for instance
\cite{Anderson_Dostert_Grabner+2023:riesz_green_energy,
  Beltran_Marzo_Ortega-Cerda2016:determinantal,
  Beltran-Etayo2018:projective_ensemble, Hirao2021:finite_frames_determinantal,
  Beltran_Ferizovic2020:approximation_uniform_SO3,
  Marzo_Ortega-Cerda2018:expected_riesz_energy,
  Hirao2018:qmc_designs_determinantal}). Since the notion of hyperuniformity
has been motivated as a property of point sets, which behave ``better than
i.i.d. random'', it is natural to test samples of determinantal processes for
this property.

The rotation invariant processes on the spaces $\mathbb{FP}^{d-1}$ have been
characterised in \cite{Anderson_Dostert_Grabner+2023:riesz_green_energy}. The
most natural amongst them is the process induced by the projection kernel on
the spaces of eigenfunctions for the $N+1$ smallest eigenvalues of the Laplace
operator. This kernel is given by
\begin{equation}
  \label{eq:harmonic-kernel}
  \begin{split}
    K_N(x,y)&=\sum_{n=0}^N\sum_{k=1}^{m_n}Y_{n,k}(x)Y_{n,k}(y)=
    \sum_{n=0}^N\frac{m_n}{P_n^{(\alpha,\beta)}(1)}
    P_n^{(\alpha,\beta)}(\cos(2\theta(x,y)))\\
    &=
    \frac{(\alpha+\beta+2)_N}{(\beta+1)_N}P_N^{(\alpha+1,\beta)}(\cos(2\theta(x,y)))
  \end{split}
\end{equation}
(see, for instance \cite{Andrews_Askey_Roy1999:special_functions}). This
process produces
\begin{equation*}
  K_N(x,x)=\frac{(\alpha+\beta+2)_N(\alpha+2)_N}{N!(\beta+1)_N}\sim
  \frac{\Gamma(\beta+1)}{\Gamma(\alpha+\beta+2)\Gamma(\alpha+2)}N^{2\alpha+2}
\end{equation*}
points; we use the notation $\mathscr{X}_N^H$ for this process and call it the
\emph{harmonic process} like the corresponding process on the sphere introduced
in \cite{Beltran_Marzo_Ortega-Cerda2016:determinantal}. For more details
on determinantal point processes we refer to
\cite{Hough_Krishnapur_Peres+2009:zeros_gaussian} again.

We notice that the number variance of $\mathscr{X}_N^H$ equals the expectation
of the sum \eqref{eq:v=gr} under law of the process $\mathscr{X}_N^H$,
\begin{equation}
  \label{eq:var-harmonic}
  \mathbb{V}(\mathscr{X}_N^H,r)=\mathbb{E}\sum_{x,y}g_r(x,y).
\end{equation}
Such expectations can be expressed easily by the general theory of
determinantal processes
\begin{multline*}
  \mathbb{E}\sum_{x,y}g_r(x,y)=\mathbb{E}\sum_{x\neq y}g_r(x,y)+
  \mathbb{E}\sum_{x}g_r(x,x)\\
  =\iint\limits_{(\mathbb{FP}^{d-1})^2}
  g_r(x,y)\left(K_N(x,x)K_N(y,y)-K_N(x,y)K_N(y,x)\right)\,d\sigma(x)\,d\sigma(y)
  \\+K_N(\cdot,\cdot)g_r(\cdot,\cdot).
\end{multline*}
Since all functions occurring in this last equation are zonal, we can use
\eqref{eq:zonal} to further simplify this equation
\begin{multline*}
  \mathbb{V}(\mathscr{X}_N^H,r)=\\C_{\alpha,\beta}
  \int_0^{\frac\pi2}g_r(\cos(2\theta))
  \left(K_N(1)^2-K_N(\cos(2\theta))^2\right)
  \sin(\theta)^{2\alpha+1}\cos(\theta)^{2\beta+1}\,d\theta\\+K_N(1)g_r(1);
\end{multline*}
here we have made use of the convention to write $f(\cos(2\theta(x,y)))=f(x,y)$
for zonal functions.
Using the facts that
\begin{align*}
  &C_{\alpha,\beta}\int_0^{\frac\pi2}K_N(\cos(2\theta))^2
  \sin(\theta)^{2\alpha+1}\cos(\theta)^{2\beta+1}\,d\theta=K_N(1)\\
  &C_{\alpha,\beta}\int_0^{\frac\pi2}g_r(\cos(2\theta))
  \sin(\theta)^{2\alpha+1}\cos(\theta)^{2\beta+1}\,d\theta=0
\end{align*}
we obtain
\begin{multline*}
  \mathbb{V}(\mathscr{X}_N^H,r)=\\C_{\alpha,\beta}\int_0^{\frac\pi2}
  \left(g_r(1)-g_r(\cos(2\theta))\right)K_N(\cos(2\theta))^2
  \sin(\theta)^{2\alpha+1}\cos(\theta)^{2\beta+1}\,d\theta.
\end{multline*}

We now remark that
\begin{multline*}
  g_r(1)-g_r(\cos(2\theta(x,y)))=\frac12\sigma(B(x,r)\triangle B(y,r))\\\leq
  C'\theta(x,y)\mathrm{surface}(\partial B(\cdot,r))\leq C''\sin(\theta)r^{D-1}
\end{multline*}
using \eqref{eq:symmetric}. Thus we have
\begin{equation*}
  \mathbb{V}(\mathscr{X}_N^H,r)=\mathcal{O}\left(r^{D-1}\int_0^{\frac\pi2}
  K_N(\cos(2\theta))^2
  \sin(\theta)^{2\alpha+2}\cos(\theta)^{2\beta+1}\,d\theta\right),
\end{equation*}
and it remains to estimate the integral
\begin{equation*}
  \frac{(\alpha+\beta+2)_N^2}{(\beta+1)_N^2}\int_0^{\frac\pi2}
  P_N^{(\alpha+1,\beta)}(\cos(2\theta))^2\sin(\theta)^{2\alpha+2}
  \cos(\theta)^{2\beta+1}\,d\theta;
\end{equation*}
notice that this is an integral over the square of a Jacobi polynomial against
the ``wrong'' weight function. We use the estimate \eqref{eq:jacobi-bound} for
$a=\alpha+1$ and $b=\beta$ in
the interval $[\frac cN,\frac\pi2]$ and the trivial bound
$|P_N^{(\alpha+1,\beta)}(\cos(2\theta))|\leq P_N^{(\alpha+1,\beta)}(1)$ in the
interval $[0,\frac CN]$ to obtain
\begin{multline*}
  \int_0^{\frac\pi2}
  P_N^{(\alpha+1,\beta)}(\cos(2\theta))^2\sin(\theta)^{2\alpha+2}
  \cos(\theta)^{2\beta+1}\,d\theta\\=
  \mathcal{O}\left(\int_0^{\frac CN}N^{2\alpha+2}\theta^{2\alpha+2}
    \,d\theta\right)+
\mathcal{O}\left(\int_{\frac
    CN}^{\frac\pi2}\frac1{N\sin(\theta)}\,d\theta\right)\\=
\mathcal{O}\left(\frac1N\right)+\mathcal{O}\left(\frac{\log(N)}N\right)=
\mathcal{O}\left(\frac{\log(N)}N\right).
\end{multline*}
Summing up, we have obtained
\begin{equation*}
  \mathbb{V}(\mathscr{X}_N^H,r)=
  \mathcal{O}\left(r^{D-1}\frac{(\alpha+\beta+2)_N^2}{(\beta+1)_N^2}
    \frac{\log(N)}N\right)=\mathcal{O}\left(r^{D-1}\frac{K_N(1)\log(N)}N\right).
\end{equation*}
Thus we have shown the following theorem.
\begin{theorem}
  The harmonic process $\mathscr{X}_N^H$ is hyperuniform for large and small
  balls. For balls at threshold order the weaker relation
  \begin{equation*}
    \mathbb{V}(\mathscr{X}_N^H,rK_N(1)^{-\frac1D})
    =\mathcal{O}\left(r^{D-1}\log(N)\right)
  \end{equation*}
  holds.
\end{theorem}
%%% Local Variables:
%%% mode: latex
%%% TeX-master: "Main"
%%% End:
